\chapter{Kết luận và Hướng phát triển}
\label{Chapter4}

\section{Tổng kết công trình nghiên cứu}

Thông qua quá trình thực hiện, đề tài đã triển khai hệ thống học từ vựng trực tuyến, tích hợp các nguyên lý khoa học nhận thức và cơ chế trò chơi hóa trong một nền tảng web. Về mặt kỹ thuật, hệ thống vận hành ổn định với 231/231 bài kiểm thử đơn vị thành công và thời gian phản hồi P95 đạt 520ms dưới tải 100 người dùng đồng thời. Cần phân biệt rõ rằng các kết quả này chứng minh \textit{tính đúng đắn kỹ thuật} của hệ thống — rằng các tính năng hoạt động đúng theo đặc tả — chứ chưa phải bằng chứng sư phạm về hiệu quả cải thiện ghi nhớ hay duy trì động lực học tập dài hạn. Dự án đã hiện thực hóa một nền tảng có tiềm năng kiểm nghiệm các giả thuyết về Gamification trong giáo dục ngôn ngữ, và việc kiểm nghiệm đó là hướng nghiên cứu cần tiếp nối trong tương lai.

Dưới góc độ chuyên môn, dự án đã hệ thống hóa thành công thuật toán SM-2 để thiết lập một lộ trình ôn tập cá nhân hóa, dựa trên sự phân tích lịch sử phản hồi của từng chủ thể. Việc đa dạng hóa các phương thức tương tác từ Flashcard đến rèn luyện thính giác đã tạo ra một môi trường thực hành truy hồi (Active Recall) đa chiều. Cách tiếp cận này không chỉ giúp củng cố các dấu vết trí nhớ từ nhiều giác quan mà còn gia tăng hiệu quả của việc chuyển dịch thông tin từ trí nhớ ngắn hạn sang lưu trữ dài hạn một cách bền vững.

Điểm thiết kế đặc trưng của Gamification nằm ở việc chuyển dịch trọng tâm từ các áp lực thành tích ngoại cảnh sang cơ chế nuôi dưỡng động lực nội sinh. Thay vì chỉ sử dụng các chỉ số bề nổi, hệ thống đã xây dựng cơ chế gắn kết thông qua thú ảo. Sự phát triển của thú ảo gắn trực tiếp với tiến trình học tập, tạo ra phản hồi tích lũy nhằm đáp ứng nhu cầu về năng lực (Competence) và tự chủ (Autonomy) theo SDT. Việc gắn kết tiến trình tich lũy kiến thức với sự phát triển của thú ảo giúp giảm bớt tính đơn điệu thường gặp trong các công cụ học từ vựng hiện có.

Về phương diện kiến trúc, việc vận hành Onion Architecture trên hệ sinh thái .NET 8 và ReactJS đã đảm bảo cho hệ thống một nền tảng vững chắc và khả năng bảo trì tối ưu. Việc phân tách rạch ròi giữa logic nghiệp vụ cốt lõi và hạ tầng kỹ thuật giúp ứng dụng linh hoạt trước các thay đổi công nghệ. Sự phối hợp của các mẫu thiết kế hiện đại cùng giao thức SignalR đã tạo ra một hạ tầng có độ trễ thấp, sẵn sàng đáp ứng các yêu cầu mở rộng quy mô phục vụ cộng đồng người dùng lớn trong tương lai.

\section{Những hạn chế và rào cản tồn tại}

Mặc dù đạt được những bước tiến quan trọng về kỹ thuật, dự án vẫn còn những khía cạnh cần được nhìn nhận khách quan để cải thiện. Trước hết, cơ chế đối kháng (PvP) hiện tại mới dừng lại ở việc thiết lập các quy tắc logic cơ bản. Việc cân bằng các chỉ số chiến đấu và nâng cấp hiệu ứng hình ảnh để tối ưu hóa trải nghiệm thời gian thực vẫn là một thách thức kỹ thuật lớn, đòi hỏi các giải pháp xử lý băng thông và thuật toán tại Server chuyên sâu hơn.

Bên cạnh đó, khả năng tương thích đa thiết bị vẫn là một giới hạn đáng kể. Các giao diện quản trị chứa mật độ dữ liệu lớn chưa đạt đến độ tối ưu hoàn hảo trên các màn hình di động nhỏ, gây khó khăn cho việc học tập mọi lúc mọi nơi. Đồng thời, do giới hạn về nguồn lực công nghệ tích hợp, hệ thống hiện tại chủ yếu tập trung vào kỹ năng đọc hiểu và ghi nhớ ngữ nghĩa, chưa thể triển khai các phân hệ nhận diện giọng nói (Speech-to-Text) để hỗ trợ toàn diện kỹ năng phát âm và giao tiếp chủ động cho người dùng.

Hạn chế quan trọng nhất của công trình, cần được nhìn nhận thẳng thắn, là sự vắng mặt hoàn toàn của nghiên cứu người dùng (user study) có kiểm soát. Trong khi đề tài đặt ra bài toán nghiên cứu về hiệu quả học tập và động lực dài hạn — dẫn chứng các lý thuyết SDT, đường cong Ebbinghaus và thuật toán SM-2 — toàn bộ kết quả trình bày chỉ xác nhận hệ thống \textit{vận hành đúng theo đặc tả kỹ thuật}, chứ chưa phải bằng chứng rằng hệ thống \textit{cải thiện khả năng ghi nhớ hay duy trì động lực} so với các phương pháp hiện có. Các tuyên bố liên quan đến Gamification — rằng cơ chế thú ảo hỗ trợ SDT hay thúc đẩy sự gắn kết dài hạn — ở giai đoạn này vẫn là giả thuyết thiết kế có căn cứ lý thuyết, chứ chưa phải kết luận được kiểm chứng qua dữ liệu người dùng thực. Để có thể đưa ra những khẳng định mang tính sư phạm, công trình cần được tiếp nối bằng một giai đoạn thực nghiệm có đối chứng, đo lường ít nhất: (1) tỉ lệ duy trì người dùng sau 4--8 tuần so với nhóm đối chứng sử dụng phương pháp học thẻ ghi nhớ thông thường; (2) hiệu quả ghi nhớ dài hạn qua kiểm tra trì hoãn (delayed retention test); (3) mức độ tự báo cáo động lực học tập theo thang đo chuẩn hóa \cite{RyanDeci2017}. Đây là khoảng trống dữ liệu quan trọng nhất cần ưu tiên khắc phục trong các giai đoạn nghiên cứu tiếp theo.

\section{Tầm nhìn và lộ trình hoàn thiện}

Hướng tới mục tiêu chuyển đổi từ một công trình nghiên cứu sang một giải pháp EdTech hoàn chỉnh, lộ trình tương lai sẽ tập trung vào việc chuyên sâu hóa các thành tố Gamification. Việc xây dựng hệ thống kỹ năng đa tầng và chiến thuật đối kháng phức hợp cho thú ảo sẽ biến các phiên học tập thành những trận đấu trí tuệ, đòi hỏi sự kết hợp giữa kiến thức ngôn ngữ và tư duy logic.

Trên phương diện công nghệ, việc ứng dụng Trí tuệ nhân tạo (AI) sẽ là ưu tiên hàng đầu để phân tích sâu sắc hành vi học tập và đặc tính quên lãng của từng cá nhân, từ đó tự động hóa việc kiến tạo lộ trình ôn tập tối ưu. Sự kết hợp của công nghệ xử lý ngôn ngữ tự nhiên (NLP) và nhận diện giọng nói sẽ mở rộng hệ sinh thái sang các kỹ năng nói và phát âm, tạo ra một môi trường rèn luyện ngôn ngữ toàn diện.

Về mặt quy mô cộng đồng, hệ thống định hướng xây dựng một quần thể học thuật bền vững thông qua các tính năng bang hội và sự kiện liên đấu. Để gia tăng tính tiện dụng, việc phát triển ứng dụng di động chính thức trên các nền tảng React Native hoặc Flutter sẽ được triển khai. Toàn bộ hệ thống sẽ được vận hành trên hạ tầng điện toán đám mây chuyên nghiệp với các cơ chế bảo mật và cân bằng tải tự động, nhằm phát triển hệ thống theo hướng tích hợp học tập và yếu tố giải trí.